\documentclass[conference]{IEEEtran}
\usepackage[T1]{fontenc}
\usepackage[utf8]{inputenc}
\usepackage{amsmath,amssymb}
\usepackage{booktabs}
\usepackage{array}
\usepackage{url}
\usepackage[hidelinks]{hyperref}
\newcommand{\model}[1]{\texttt{#1}}
\newcommand{\checkc}[1]{\textsc{#1}}

\title{Auditing a Conflicted Judge: Hybrid Verification for an\\
LLM-Scored Forensic Reliability Benchmark}
\author{Anonymous submission}

\begin{document}
\maketitle

\begin{abstract}
Generative-model benchmarks increasingly use another generative model as the
judge, and that judge may be related to the systems it scores. We study this
problem using Android forensics as a case study. Thirty synthetic cases and six systems yield 180 evidence-grounded
reports. A five-dimension claim rubric measures support, provenance,
observation--inference separation, attribution, and reproducibility, but its
primary rater belongs to the same model family as three evaluated systems. A
cross-family diagnostic does not resolve the conflict: pooled Krippendorff's
$\alpha$ is 0.052, although differing rating units confound that estimate. A
direct same-unit test is more informative: on 60 claims scored blind by two
authors, the primary judge exceeds the authors by $+0.20$ points on
Claude-authored claims and by $-0.11$ on Qwen-authored claims (difference
0.31, response-clustered permutation $p=0.003$), and the gap persists between
the two compact systems ($p=0.011$). We
therefore introduce a hybrid benchmark architecture: interpretive constructs
remain rubric-rated, while machine-decidable properties are checked against
structured case facts. Four versioned checks flag nine responses across the six
systems. Two model-blinded author raters separately confirm seven flags and mark
two as unclear. In a stratified sample of 24 nonflagged responses, they agree on
one missed violation and 22 nonviolations, and disagree on one additional C3
case; this sample does
not estimate full-system recall. A conformance suite containing 20 seeded violations
and 20 matched non-violations per check detects all 80 positives and rejects all
80 negatives. These fixtures validate implementation behavior rather than
natural prevalence or recall. The resulting benchmark reports which conclusions
depend on the conflicted judge and which can be recomputed without it.
\end{abstract}

\section{Introduction}
When one generative model judges another, the resulting score may combine
output quality, evaluator preference, presentation style, and family affinity.
A benchmark built on such scores can look precise while its measurement is
affected by the same kind of model behavior it is meant to evaluate. Treating
the scores as ground truth can make an uncertain measurement look like a stable
leaderboard.

Forensic reporting is a useful setting for studying this problem. An evidentiary
statement can be fluent and still be unreliable. A forensic assistant may
accurately report that a domain appears in an application's strings, then
overstate that the application contacted that domain. It may describe activity
observed on a device as an action by a particular person. In both examples, the
failure concerns the relationship between a claim and its evidence, which
lexical similarity to a reference answer does not measure.

Claim-level evaluation offers a better starting point. Long responses can be
decomposed into atomic findings and each finding scored for support and
traceability~\cite{factscore,safe}. Yet this approach creates a measurement
problem when another LLM assigns the scores. Model judges exhibit self-preference
and presentation effects~\cite{panickssery2024,zheng2023}. If the judge belongs
to the same family as systems under test, a high score can reflect system
quality, judge affinity, or both.

We examine the broader measurement problem through a controlled
Android-forensics benchmark. Thirty synthetic cases combine manifest entries,
extracted strings, certificate metadata, network indicators, and traffic logs.
Six systems from three providers produce one report per case. Four systems also
have claim-level scores from a primary LLM rater. We treat those scores as
measurements produced by a conflicted instrument, then ask which conclusions
require interpretation and which can be decided directly from structured facts.

This paper contributes: (1) a hybrid benchmark architecture that assigns
interpretive and machine-decidable properties to different measurement
instruments; (2) a 30-case benchmark and five-dimension rubric for forensic
report reliability; (3) a direct, family-stratified audit of the conflicted
judge against blinded author scores on identical claims, which finds leniency
concentrated on same-family outputs; (4) a deterministic verifier with four inspectable checks, witnesses,
  and a 160-fixture conformance suite; and (5) a six-system audit spanning
  Anthropic, OpenAI, and a local Qwen model, supplemented by a blinded two-author
  audit. Deterministic checks support specific
statements about explicit failure modes. They neither validate the full rubric
nor convert a conflicted score table into ground truth.

We organize the study around four questions. \textbf{RQ1:} What result does the
original model-rated instrument produce, and how sensitive is that measurement
to a cross-family diagnostic? \textbf{RQ2:} Which pre-specified reliability
failures can be decided without a generative rater? \textbf{RQ3:} Does the
executable verifier behave correctly on seeded positive and matched negative
controls? \textbf{RQ4:} What does a model-blinded two-author audit reveal about
the rubric ratings, the judge's family affinity, and verifier coverage? Keeping
these questions separate avoids merging agreement statistics, detector behavior,
and natural model errors into a single accuracy figure.

\section{Related Work}
\subsection{Forensic evaluation}
Studiawan \emph{et al.} study standardized evaluation for LLM-assisted forensic
timelines and include BLEU and ROUGE~\cite{studiawan2025}. AutoDFBench evaluates
generated forensic code~\cite{autodfbench}. Our target is whether each
natural-language finding remains within what the supplied artifacts justify.

FActScore introduced atomic factual precision~\cite{factscore}. SAFE uses a
search-augmented model evaluator~\cite{safe}, while ALCE and attribution-scoring
work evaluate whether cited evidence supports generated text~\cite{alce,attrscore}.
For forensic reporting, source support must be combined with provenance,
epistemic status, attribution restraint, and reconstructable reasoning.

\subsection{Model judges and benchmark validity}
Model-based judges scale beyond human panels but inherit systematic biases.
They may favor their own generations~\cite{panickssery2024}, and outcomes can
change with answer order or style~\cite{zheng2023,faireval}. G-Eval and
Prometheus show the promise of rubric-guided model evaluation, including strong
correlation with human ratings in their studied settings~\cite{geval,prometheus},
but neither result makes a new judge ground truth for a domain-specific
evidentiary task. Agreement coefficients are useful only when raters measure the
same units and constructs; values below 0.667 do not support even tentative
conclusions under commonly used guidance~\cite{krippendorff}. We respond by
separating model-rated measurements from properties that executable checks can
decide.

\subsection{Positioning and novelty}
Prior work supplies three adjacent components: benchmarks for forensic tasks,
research on bias and instability in model judges, and factuality or attribution
methods that decompose long answers into supported claims. These components do
not by themselves specify how to publish a forensic benchmark whose rating
instrument is acknowledged to be conflicted. Our contribution is an explicit
measurement boundary. The benchmark declares which constructs need
interpretation, compiles selected case facts into executable conformance
properties, publishes witnesses and applicability sets, and reports the two
evidence streams separately. The design remains useful even when no model judge
is demonstrably unbiased.
We do not claim that executable checks are new: unit tests in code generation,
CheckList-style behavioral tests~\cite{checklist}, and citation checkers already
apply them where ground truth permits. Our claim is narrower. We apply them to
\emph{audit} a specific conflicted judge, and we measure that judge's family
affinity directly on identical claims. We then report which results hold
without the judge.

\section{Hybrid Measurement Architecture}
The architecture separates a benchmark into an \emph{interpretive layer} and
an \emph{executable layer}. The interpretive layer covers constructs such as
whether reasoning is sufficiently traceable or whether uncertainty is expressed
appropriately. Human or model raters can apply anchored rubrics to these
constructs, and their agreement must be measured. The executable layer covers
properties that structured ground truth decides, such as whether a named domain
appears in any supplied artifact or whether an asserted network contact appears
in the traffic log. These properties are evaluated by deterministic code.

The benchmark as a whole still depends on raters, but a defined subset of its
conclusions does not, and the benchmark identifies that subset. Table
\ref{tab:architecture} maps each measurement target to its instrument, evidence,
and principal unresolved risk.

\begin{table*}[t]
\centering
\caption{Measurement boundary in the hybrid benchmark.}
\label{tab:architecture}
\small
\begin{tabular}{p{0.18\textwidth}p{0.18\textwidth}p{0.27\textwidth}p{0.25\textwidth}}
\toprule
Target & Instrument & Evidence produced & Principal limitation \\
\midrule
Support, provenance, epistemic status, attribution, reproducibility & Anchored claim rubric & Per-claim ordinal score and rationale & Rater bias, construct interpretation, and claim segmentation \\
Explicit domain contact, person attribution, configured cross-artifact relation, absent domain & Deterministic checks C1--C4 & Applicability decision, violation bit, and textual witness & Narrow coverage and unknown natural-output recall \\
Verifier implementation behavior & Seeded conformance suite & Expected and observed result for positive and matched negative fixtures & Synthetic controls do not estimate prevalence \\
Overall evidentiary quality & Human examiner & Contextual judgment over claims, response, and case & Cost, expertise, and residual disagreement \\
\bottomrule
\end{tabular}
\end{table*}

Four design rules follow. First, every executable check declares an
applicability set; a non-applicable case cannot be counted as a success. Second,
every positive decision stores a witness so a reviewer can reproduce it. Third,
zero-hit checks must be exercised by seeded positive and negative controls.
Fourth, aggregate scores must not silently combine rubric ratings with
executable violations, because the instruments answer different questions.
This use of controlled expected behavior is related to behavioral testing in
CheckList~\cite{checklist}; here, expectations are compiled from each forensic
case's structured evidence and applied to complete generated reports.

\section{Benchmark}
\subsection{Cases and task}
The benchmark contains 30 version-controlled JSON cases: 12 baseline cases,
eight attribution traps, six planted contradiction or non-corroboration cases,
and four ambiguous-timestamp cases. Each case supplies an Android manifest,
extracted strings, certificate metadata, network indicators, a traffic log, and
structured ground-truth fields. Domains and IP addresses use reserved
documentation ranges. No malware sample, personal record, or operational
infrastructure is included.

Every system receives the same single-turn instruction and the case JSON
verbatim. The requested output is a short forensic assessment that distinguishes
observations from inferences, cites the relevant artifact and field, avoids
person-level attribution without evidence, and states what additional evidence
would be required. The prompt does not expose ground truth.

Android signing semantics require care. Long-lived and self-signed application
signing certificates are common and are distinct from TLS server certificates;
Android recommends sufficiently long key validity for the expected lifetime of
an application~\cite{androidsigning}. The benchmark therefore does not treat
those properties alone as proof of maliciousness, compromise, or a broken
network trust chain.

\subsection{Case construction and controls}
Cases were designed around errors in evidentiary reasoning rather than malware
family recognition. A case can contain individually suspicious artifacts while
withholding the link needed for a stronger conclusion. A manifest permission
establishes capability, a string establishes static presence, and a traffic
record establishes an observed connection; none is substituted for another.
Attribution traps include device activity without evidence identifying the
operator. Contradiction cases place two true fields in tension and require the
report to state their relationship rather than list each field in isolation.

Each JSON object separates the scenario description, artifacts, and structured
ground-truth fields. The response harness exposes only the scenario and
artifacts. Reserved domains and addresses prevent accidental contact with live
infrastructure and make exact-domain checks reproducible. Case identifiers,
prompts, and responses are hashed so later analysis can detect accidental
changes. The synthetic design permits controlled counterfactual relations and
artifact-level disclosure without personal or operational data. Its external
validity is correspondingly narrower than a corpus of real case files.

The case distribution was fixed before adding the two OpenAI endpoints. Those
endpoints received the same serialized case objects and task instruction as the
original four systems. No model-specific prompt tuning was performed. We retain
complete outputs rather than extracting only matched phrases, allowing every
flag and rating to be traced to the original report.

\subsection{Reliability rubric}
The primary instrument assigns 0, 1, or 2 to each atomic finding along five
dimensions. \emph{Evidence support} asks whether an artifact directly supports
the finding. \emph{Provenance} asks whether the response identifies a specific
artifact and field. \emph{Observation--inference separation} asks whether the
response marks epistemic status correctly. \emph{Attribution} penalizes an
unsupported action attributed to a person or entity. \emph{Reproducibility}
asks whether another examiner could reconstruct the reasoning from supplied
artifacts. Table~\ref{tab:rubric} gives the full anchors.

\begin{table*}[t]
\centering
\caption{Claim-level forensic reliability rubric.}
\label{tab:rubric}
\small
\begin{tabular}{p{0.16\textwidth}p{0.25\textwidth}p{0.52\textwidth}}
\toprule
Dimension & Question & Score anchors \\
\midrule
Evidence support & Is the finding supported by a supplied artifact? & 0: unsupported or fabricated; 1: plausible but indirect; 2: directly supported.\\
Provenance & Can the finding be mapped to its source? & 0: no traceable source; 1: artifact type only; 2: specific artifact and field/value.\\
Observation/inference & Is epistemic status explicit and correct? & 0: inference stated as fact; 1: ambiguous; 2: observation and inference are clearly separated.\\
Attribution & Is an attributed act justified? & 0: unjustified person/entity attribution; 1: weak justification; 2: justified, qualified, or no attribution claim.\\
Reproducibility & Can another examiner reconstruct the reasoning? & 0: reasoning absent; 1: partly reconstructable; 2: reconstructable from supplied artifacts.\\
\bottomrule
\end{tabular}
\end{table*}

\subsection{Systems}
Table~\ref{tab:systems} lists the six evaluated endpoints. Each produced all 30
responses. The local Qwen model is a 7.6B Q4\_K\_M build served by Ollama. For
the two OpenAI endpoints, the harness records the requested and served snapshot,
reasoning setting, output limit, usage, case hash, prompt hash, and request
fingerprint. All 60 requests were served by the requested snapshot. The Opus
endpoint is labelled as a mixture because provider routing served 18 requests
with \model{claude-opus-4-8} and 12 with \model{claude-opus-5}; it cannot
support a version-specific conclusion.

\begin{table*}[t]
\centering
\caption{Evaluated systems. Every endpoint completed 30 cases.}
\label{tab:systems}
\scriptsize
\begin{tabular}{lll}
\toprule
Reported name & Provider/runtime & Scale group \\
\midrule
Claude Opus endpoint & Anthropic API & Frontier \\
Claude Sonnet 5 & Anthropic API & Frontier \\
Claude Haiku 4.5 & Anthropic API & Compact \\
Qwen2.5 7B & Local Ollama & Compact \\
GPT-5.5 (2026-04-23) & OpenAI API & Frontier \\
GPT-5.4 Mini (2026-03-17) & OpenAI API & Compact \\
\bottomrule
\end{tabular}
\end{table*}

\section{Evaluation Design}
\subsection{Primary and cross-family raters}
The primary rater, Claude Sonnet 5, decomposed responses from the four original
endpoints into findings and scored all five dimensions with short rationales.
This produced 830 annotated claims from 120 responses. Model identity was
visible. Because the rater is identical to one evaluated system and belongs to
the same family as two others, these values describe the instrument's output;
they are not independent ground truth.

A local Qwen2.5 7B rater blindly re-scored a fixed stratified sample of 30
Claude-authored responses, ten per Claude endpoint. It assigned one score per
dimension to the whole response. We compare it with rounded response-level
means of primary claim scores using ordinal Krippendorff's $\alpha$. This is a
diagnostic because the units differ: primary scores are averaged over findings,
while secondary scores are holistic. The sample contains 21 unique cases, so we
also aggregate by case before a paired Wilcoxon test.

\subsection{Deterministic verifier}
The verifier operates on response text and structured case facts. It emits at
most one response-level violation per check and stores a witness sentence or
the missing fact pair. The checks are: \checkc{C1}, assertion of contact with a
domain found only in static strings and absent from traffic; \checkc{C2}, a
concrete person-action attribution without qualification; \checkc{C3}, failure
on one of six configured cases to mention and relate two facts that form a
contradiction or non-corroboration pattern; and \checkc{C4}, a reserved domain
absent from every supplied artifact.

The domain parser is restricted to the benchmark's reserved
\texttt{example.com} namespace so Java package names and filenames are not
mistaken for network indicators. Negation, capability statements, and
prospective recommendations are excluded from \checkc{C1}. Applicability per
model is 4 cases for \checkc{C1}, 30 for \checkc{C2}, 6 for \checkc{C3}, and 30
for \checkc{C4}.

For a response $r$ and case $x$, let $D_r$ be reserved domains extracted from
assertive response sentences, $D_s(x)$ domains occurring in static artifacts,
and $D_t(x)$ domains in observed traffic. \checkc{C1} fires when an assertion
maps to $d\in D_s(x)\setminus D_t(x)$ and does not fall under the negation,
capability, or recommendation exclusions. \checkc{C4} instead fires for
$d\in D_r\setminus D_a(x)$, where $D_a(x)$ is the union of all case-artifact
domains. \checkc{C2} searches for a concrete person-action construction and
then rejects qualified or explicitly uncertain language. \checkc{C3} is
applicable only to six cases. It requires lexical evidence for each member of a
configured fact pair and relation language that marks contradiction,
non-corroboration, or temporal inconsistency. The implementation returns the
triggering sentence for C1, C2, and C4, and the missing relation for C3.

\begin{table*}[t]
\centering
\caption{Executable properties and their scope. Counts are per evaluated system.}
\label{tab:checks}
\small
\begin{tabular}{p{0.06\textwidth}p{0.38\textwidth}p{0.10\textwidth}p{0.33\textwidth}}
\toprule
Check & Violation property & Applicable cases & Main safeguards \\
\midrule
C1 & Asserts communication with a domain present only in static strings & 4 & Exact reserved domain; traffic absence; negation, capability, and prospective-language exclusions \\
C2 & Attributes a concrete action to a person without evidence or qualification & 30 & Action construction plus uncertainty and qualification exclusions \\
C3 & Omits the configured relation between two contradiction or non-corroboration facts & 6 & Both facts and accepted relation terms are case-configured \\
C4 & Names a reserved domain absent from every supplied artifact & 30 & Exact reserved-domain extraction and artifact-wide set comparison \\
\bottomrule
\end{tabular}
\end{table*}

Every verifier flag and sampled nonflag was subsequently reviewed against the
complete response and case evidence under the model-blinded author-audit
protocol below. The verifier was specified after the initial model-scored
analysis; its results are an audit on this dataset rather than performance of a
preregistered detector.

\subsection{Seeded conformance suite}
Natural outputs provide no positive examples for C2 or C4. A zero count could
therefore reflect robust model behavior or a detector that never fires. We test
implementation behavior using 20 seeded positive and 20 matched negative
fixtures for each check. C1 controls vary affirmative contact claims against
negation, capability, and future recommendations. C2 contrasts unsupported
person-action statements with qualified interpretations. C3 pairs reports that
state both facts and their required relation with reports that omit the
relation. C4 inserts absent reserved domains and matches them with domains
present in the artifacts.

The 160 fixtures call the same verifier functions used for model outputs; they
are not a separate reimplementation. A fixture passes only when its expected
violation set exactly matches the observed set. These controls test parser and
rule behavior at known boundaries. Because the language is deliberately seeded,
they do not measure recall over naturally occurring paraphrases, error
prevalence, or generalization to new case schemas.

\subsection{Model-blinded author-audit protocol}
We prepared a fixed packet for two authors with security and ML-evaluation
backgrounds. The authors report completing it separately without seeing model
identity, primary scores, the coordinator key, or each other's review. The
packet samples 60 atomic claims, 15 from each of the four systems with primary
claim ratings; presents all nine verifier flags; and samples 24 nonflagged
responses, four per system. For every item it includes the case artifacts and
complete response. Reviewers apply the five rubric dimensions, judge each flag
as valid, invalid, or unclear, and inspect nonflags for missed C1--C4 violations.

We report exact agreement and quadratic-weighted Cohen's $\kappa$ for each
ordinal dimension, agreement with the primary judge, and exact agreement for
flag and miss decisions. We retain unclear decisions without adjudication. The
reviews were completed by authors rather than external examiners. The authors
report that all scores and decisions were selected and filled manually. They
consulted ordinary websites and reference material; any LLM use was limited to
editorial phrasing after the judgments. We therefore treat these results as an
author audit, not independent human validation.

\emph{Family-stratified judge audit.} Because the 60 claims are balanced over
the four primary-scored systems and scored on identical units by the judge and
both authors, they permit a direct affinity test. For each claim we define
\emph{judge excess} as the primary score minus the mean author score, averaged
over the five dimensions. A lenient but unbiased judge would show similar
excess across systems. We compare Claude-authored claims with Qwen-authored
claims, the judge's own endpoint (Sonnet) with all others, and the two compact
systems (Haiku and Qwen), which holds scale approximately fixed. Inference uses
a two-sided permutation test with 20{,}000 permutations. System labels are
permuted at the response level (56 responses) so that claims from one response
move together. This analysis was specified after the audit data were collected
and is therefore post hoc.

\subsection{Statistical analysis}
For the primary rater, we average within case and then across 30 cases. Pairwise
differences use paired Wilcoxon tests with Holm correction. Cross-rater
agreement uses ordinal $\alpha$ and paired Wilcoxon tests at response and
case-aggregated levels. For deterministic checks, we report counts rather than
an aggregate accuracy. An exploratory comparison groups the three frontier
endpoints and the three compact endpoints. Each endpoint has ten applicable
\checkc{C1}+\checkc{C3} case-checks, giving 30 observations per group, and we
apply a two-sided Fisher exact test. Repeated cases and heterogeneous checks
limit this test; it is supporting evidence only.

\section{Results}
\subsection{The model-rated result}
Table~\ref{tab:judge} shows case-weighted means from the 830 claim annotations.
The primary rater assigns near-ceiling values to the Opus endpoint and Sonnet,
and lower values to Haiku and Qwen2.5. Paired tests separate Opus and Sonnet
from Haiku and Qwen2.5 after Holm correction
($p_{\mathrm{Holm}}\leq1.5\times10^{-5}$). Opus and Sonnet are not separated
($p_{\mathrm{Holm}}=0.655$), nor are Haiku and Qwen2.5
($p_{\mathrm{Holm}}=0.369$). These comparisons remain judge-dependent.

\begin{table}[t]
\centering
\caption{Primary-rater means (0--2), averaged within case and then across 30 cases.}
\label{tab:judge}
\scriptsize
\begin{tabular}{lccccc}
\toprule
Model & Supp. & Prov. & O/I & Attr. & Repro.\\
\midrule
Opus endpoint & 1.988 & 1.996 & 2.000 & 2.000 & 1.996\\
Claude Sonnet & 1.985 & 2.000 & 2.000 & 2.000 & 2.000\\
Claude Haiku & 1.598 & 1.892 & 1.595 & 2.000 & 1.602\\
Qwen2.5 7B & 1.758 & 1.771 & 1.790 & 2.000 & 1.728\\
\bottomrule
\end{tabular}
\end{table}

\subsection{The cross-rater diagnostic does not reproduce the primary result}
Pooled agreement over 150 dimension scores is $\alpha=0.052$. Evidence support
and reproducibility each yield $\alpha=-0.270$; observation--inference yields
$\alpha=0.356$. Provenance and attribution have insufficient variance for a
defined coefficient. The secondary rater is 0.187 points stricter on average.
The difference is significant at the response level ($p=0.00034$) and after
aggregation by unique case ($p=0.00205$). These results show that the original
measurement is not reproduced under this cross-family protocol. They do not
isolate evaluator family as the cause: changing both the rater and the unit of
analysis prevents interpreting either score set as correct.

\subsection{The author audit does not validate the primary judge}
Reviewer A and Reviewer B independently completed model-blinded manual reviews.
Each completed all 300 requested ordinal scores and all 33 response-level
decisions. Table~\ref{tab:human} shows the claim results.
Exact author--author agreement ranges from 88.3\% to 100\%; quadratic-weighted
$\kappa$ ranges from 0.798 to 0.951 where defined. Provenance agreement is
undefined because both authors assign the constant score 2. High agreement here
should not be read as independent external replication because both raters are
authors.

Against the primary judge, exact agreement ranges from 55.0\% for support to
88.3\% for attribution. Quadratic-weighted $\kappa$ ranges from 0 to 0.503.
The zero values for provenance and attribution arise when one rater is constant,
so the corresponding high exact agreement does not imply chance-corrected
agreement. The audit therefore supplies a reproducible sensitivity check, not
ground truth for the five-dimension rubric.

\begin{table*}[t]
\centering
\caption{Model-blinded author audit of 60 sampled claims. A and B are the two
authors; P is the primary judge. ``Exact'' is percent exact agreement and
$\kappa_w$ is quadratic-weighted Cohen's $\kappa$.}
\label{tab:human}
\scriptsize
\begin{tabular}{lcccccc}
\toprule
Dimension & A mean & B mean & A--B exact & A--B $\kappa_w$ & A--P exact; $\kappa_w$ & B--P exact; $\kappa_w$ \\
\midrule
Support & 1.17 & 1.22 & 88.3 & 0.798 & 55.0; 0.404 & 55.0; 0.331 \\
Provenance & 2.00 & 2.00 & 100.0 & undef. & 83.3; 0.000 & 83.3; 0.000 \\
Obs./inference & 1.55 & 1.55 & 90.0 & 0.888 & 65.0; 0.397 & 65.0; 0.397 \\
Attribution & 1.87 & 1.85 & 98.3 & 0.951 & 88.3; 0.000 & 88.3; 0.000 \\
Reproducibility & 1.58 & 1.57 & 91.7 & 0.830 & 73.3; 0.503 & 71.7; 0.477 \\
\bottomrule
\end{tabular}
\end{table*}

\subsection{The judge is more lenient toward same-family outputs}
Table~\ref{tab:family} directly tests the affinity that motivates this paper.
On the same 60 claims, the primary judge exceeds the blinded author mean by
$+0.202$ points for Claude-authored claims and by $-0.107$ for Qwen-authored
claims. The difference of 0.309 is significant under the response-clustered
permutation test ($p=0.0025$). Scale does not explain the gap. The judge gives
Haiku and Qwen nearly identical means (1.587 and 1.573), but the authors rate
Haiku lower (1.320 and 1.680). The resulting excess difference is 0.373
($p=0.011$). The judge's own endpoint does not differ detectably from the other
three systems ($+0.162$, $p=0.139$). Thus the evidence points to family-level
leniency rather than strict self-preference. By dimension, the Claude--Qwen gap
is largest for observation--inference separation ($+0.47$, $p=0.017$), followed
by support ($+0.30$, $p=0.086$) and reproducibility ($+0.19$, $p=0.23$).

This result partly attributes the leniency visible in
Table~\ref{tab:human}: the authors' mean support score of about 1.2 against the
judge's 1.58 is concentrated in same-family outputs. The result has several
limits. There are 15 claims per system. The judge controlled claim segmentation. The
authors saw no model identity but could, in principle, recognize family style.
The reference scores come from authors rather than external examiners. The
analysis is also post hoc. The test cannot separate family affinity from a
preference for a presentation style that Claude endpoints share. The self-preference literature
itself links these mechanisms~\cite{panickssery2024}. Within these limits,
family-level leniency is measured on this sample rather than assumed.

\begin{table}[t]
\centering
\caption{Judge excess on the 60 audited claims (primary score minus the mean of
the two author scores, averaged over dimensions). $p$ values are two-sided
response-clustered permutation tests.}
\label{tab:family}
\scriptsize
\begin{tabular}{lccccc}
\toprule
Source & Claims & Resp. & Primary & Authors & Excess \\
\midrule
Opus endpoint & 15 & 13 & 1.933 & 1.840 & $+0.093$ \\
Sonnet (judge) & 15 & 13 & 1.947 & 1.700 & $+0.247$ \\
Claude Haiku & 15 & 15 & 1.587 & 1.320 & $+0.267$ \\
Qwen2.5 7B & 15 & 15 & 1.573 & 1.680 & $-0.107$ \\
\midrule
Claude (pooled) & 45 & 41 & 1.822 & 1.620 & $+0.202$ \\
\midrule
\multicolumn{4}{l}{Contrast} & $p$ & Excess diff. \\
\midrule
\multicolumn{4}{l}{Claude $-$ Qwen} & 0.003 & $+0.309$ \\
\multicolumn{4}{l}{Haiku $-$ Qwen (scale-matched)} & 0.011 & $+0.373$ \\
\multicolumn{4}{l}{Sonnet $-$ other three (self)} & 0.139 & $+0.162$ \\
\bottomrule
\end{tabular}
\end{table}

\subsection{Rater-free checks across six models}
Table~\ref{tab:det} reports the deterministic audit. Nine flags occur among 180
responses. Both author raters confirm seven, mark none invalid, and mark two
\checkc{C3} flags on case AND-022 unclear because the review packet does
not expose the exact configured fact pair. We therefore do not report a 9/9
precision estimate. Zero \checkc{C2} or \checkc{C4} flags means only that the
specified detector found no matching errors.

\begin{table}[t]
\centering
\caption{Response-level verifier flags before author review. Applicability per model is C1: 4, C2: 30, C3: 6, C4: 30.}
\label{tab:det}
\scriptsize
\begin{tabular}{lccccc}
\toprule
Model & C1 & C2 & C3 & C4 & Total \\
\midrule
Opus endpoint & 0 & 0 & 0 & 0 & 0\\
Claude Sonnet & 0 & 0 & 0 & 0 & 0\\
Claude Haiku & 0 & 0 & 1 & 0 & 1\\
Qwen2.5 7B & 2 & 0 & 2 & 0 & 4\\
GPT-5.5 & 0 & 0 & 1 & 0 & 1\\
GPT-5.4 Mini & 0 & 0 & 3 & 0 & 3\\
\midrule
All models & 2 & 0 & 7 & 0 & 9\\
\bottomrule
\end{tabular}
\end{table}

The nonflag sample contains one natural-output miss. Both authors identify the
Qwen2.5 response for AND-021 as violating \checkc{C1} and \checkc{C3}: it says
the app communicates through a domain present only in static artifacts and does
not report that the traffic log names a different domain. They agree that 22 of
the remaining responses contain no missed C1--C4 violation. For one AND-022
response, Reviewer B records a C3 miss while Reviewer A records Unclear; this
item remains unresolved without adjudication. Thus the confirmed sampled
miss proportion is 1/24, or 4.2\% (exact 95\% binomial interval [0.1\%, 21.1\%]),
with one additional disputed item. The stratified sample is too small to
support a system-wide recall estimate.

\subsection{All seeded verifier controls pass}
Table~\ref{tab:fixtures} answers RQ3. Every seeded positive is detected and
every matched negative is rejected. The negative controls are important because
several exclusions reverse the expected outcome while retaining most lexical
content. For example, ``contacted'' should trigger C1 in an affirmative sentence,
while ``did not contact,'' ``could contact,'' and ``monitor for future contact''
should not. Likewise, a C4 absent-domain hallucination is paired with a sentence
naming a domain actually present in the case.

\begin{table}[t]
\centering
\caption{Conformance-suite results. ``Pos.'' is seeded violations detected;
``Neg.'' is matched non-violations correctly rejected.}
\label{tab:fixtures}
\small
\begin{tabular}{lccc}
\toprule
Check & Pos. & Neg. & Natural flags \\
\midrule
C1 & 20/20 & 20/20 & 2 \\
C2 & 20/20 & 20/20 & 0 \\
C3 & 20/20 & 20/20 & 7 \\
C4 & 20/20 & 20/20 & 0 \\
\midrule
Total & 80/80 & 80/80 & 9 \\
\bottomrule
\end{tabular}
\end{table}

This result removes one ambiguity around zero-hit checks: C2 and C4 can fire on
their seeded violation classes and reject their matched controls. It does not
show that the rules cover all natural phrasings or that the evaluated systems
made no C2/C4 errors outside those rule boundaries.

\subsection{Exploratory system-group contrast}
The frontier group has 1 verifier flag in 30 applicable \checkc{C1}+\checkc{C3}
case-checks; the compact group has 8 in 30. The two-sided Fisher result is
$p=0.0257$. We do not treat this as a headline ranking. The comparison combines
two checks, reuses cases across systems, and changes scale, provider, and serving
conditions simultaneously. It is a descriptive consistency check showing that
auditable failures are not confined to the original four-system evaluation.

\subsection{Observed failure modes}
The two \checkc{C1} flags both come from Qwen2.5. In one case the report turns a
domain present only in extracted strings into observed communication; in the
other, a heading and accompanying prose imply traffic despite the lack of a
traffic-log entry. The seven \checkc{C3} flags are omissions of a required
cross-artifact relation. Examples include reporting the \texttt{RECORD\_AUDIO}
permission without noting that no other supplied artifact corroborates recording
functionality, and reporting a traffic timestamp without contrasting it with a
later component \texttt{first\_seen} value. The individual facts may be copied
correctly while the evidentiary relationship that changes their interpretation
is lost.

The attribution results help explain why neither the rubric nor the verifier
registers attribution errors. The primary judge assigns 2 to every attribution
score, and \checkc{C2} never fires, although eight cases are attribution traps.
In the author audit, at least one author scores 7 of the 60 claims below 2 on
attribution. Five of these come from attribution-trap cases. They include ``direct evidence of malware
infrastructure'' and ``location data is being collected and transmitted'' stated
as fact. These are entity- and intent-level overattributions rather than the
concrete person--action constructions that \checkc{C2} targets, and the judge
did not penalize them. On this small sample the traps did elicit errors, but
\checkc{C2} is too narrow to detect them and the judge did not penalize them.
The evidence does not suggest that the traps are too easy.

The \checkc{C3} results also expose a configuration weakness. Both unclear flags
and the one disputed nonflag come from AND-022. Its configured fact pair and
accepted non-corroboration phrasings live in verifier code, not in the case
record, so the review packet did not expose them. Three uncertain decisions in
one case suggest that its relation is under-specified.
Section~\ref{sec:design} describes the schema change that addresses it.

The stronger GPT endpoint receives one \checkc{C3} flag and the compact GPT
endpoint receives three; one flag in each endpoint is unresolved in the author
audit. The raw verifier pattern is consistent with the frontier--compact
contrast, while the small number of applicable cases prevents
a stable within-provider claim (two-sided Fisher $p=0.582$ for 1/10 versus
3/10). The deterministic results support failure-mode analysis more strongly
than fine-grained model ranking.

\section{Discussion}
\subsection{Evidence that does not depend on the judge}
The evaluation yields three layers of evidence. First, the claim-level table is
useful for locating candidate differences but remains conditional on a
same-family judge. The family-stratified audit shows that this condition
matters: on identical claims, the judge's leniency relative to blinded authors
is concentrated in Claude-authored outputs. On the audited sample the authors
rate Qwen2.5 well above Haiku (1.680 versus 1.320). The judge does not register
this gap (1.573 versus 1.587). Comparisons across families in Table~\ref{tab:judge} should
therefore not be read as system rankings. Second, the cross-family rater exposes the instability of
that measurement and does not provide replacement ground truth. Third, the
deterministic checks produce nine reproducible flags without consulting either
rater; the subsequent audit confirms seven, leaves two unresolved, and finds
one missed response in a 24-response nonflag sample.

The third layer is strongest because each decision is tied to structured case
facts and an inspectable witness. Its scope is also the narrowest. The verifier
directly covers fabricated or overstated network contact and missed
cross-artifact relationships, which intersect evidence support and
observation--inference separation. It does not measure writing quality,
completeness, usefulness to an examiner, or all ways provenance and attribution
can fail.

\subsection{Implications for benchmark design}\label{sec:design}
A benchmark release should therefore have two parts. A rubric can retain human
or model judgment for constructs that require interpretation. Alongside it, a
versioned conformance suite should target properties that structured ground
truth can decide. Each rule should publish its applicable set, positive and
negative fixtures, evidence snippets, and version identifier. Rules that never
fire need seeded positive controls; otherwise zero findings conflate good model
behavior with weak detector recall.

The case schema should encode relations directly. Network provenance is cheap
to verify because observed and string-only domains occupy distinct fields.
Contradiction detection is harder because the verifier maps each of six cases
to two textual patterns and acceptable relation language. A stronger schema
would store \texttt{fact\_a}, \texttt{fact\_b}, the expected relation, and
permissible reporting variants as first-class data.

This architecture also changes how negative findings are interpreted. A zero
violation count has three possible explanations: the system avoided the error,
the check did not apply, or the implementation failed to recognize it. Explicit
applicability separates the second explanation, while seeded controls provide
evidence against simple implementation failure. Only a held-out natural sample
with independent review can address broader detector coverage. Reporting all
three stages prevents ``no flags'' from being presented as ``no errors.''

The same pattern transfers beyond forensics when a task has structured state.
In code generation, executable properties can cover compilation, tests, and
forbidden API calls while reviewers assess maintainability. In grounded question
answering, exact citation presence and source identifiers can be checked while
humans assess whether the cited passage entails the claim. In incident response,
timeline order and indicator membership can be executable while causal analysis
remains interpretive. The approach does not require complete formalization,
only a clear statement of which conclusions have been formalized.

\subsection{Trust boundary and operational use}
The verifier is part of the benchmark measurement system, not a replacement for
a forensic examiner. It assumes that the case JSON is complete for the property
being checked, that reserved-domain extraction is correct, and that case-specific
C3 relations are encoded faithfully. A corrupted or incomplete case can make a
deterministic decision reproducible but wrong. The artifact therefore includes
case records, rule source, witnesses, fixtures, version identifiers, and hashes
as one audit unit.

Operationally, a flag should be read as a review trigger with a reproducible
reason. It should not be converted into a claim about maliciousness, intent, or
legal responsibility. This distinction is especially important for Android
signing and permissions, where common platform behavior may look suspicious
without additional corroboration. The benchmark evaluates evidentiary restraint
under supplied facts; it does not certify an endpoint for autonomous casework.

\section{Limitations}
The model-blinded review was completed by the two authors, not external forensic
examiners. They report selecting and filling all scores and decisions manually.
They consulted ordinary websites and reference material; any LLM use was limited
to editorial phrasing after the judgments. The review files were completed
separately. Claim-score exact agreement ranges from 88.3\% to 100\%, flag
decisions agree 9/9, and miss decisions agree 23/24. This audit may share
assumptions with the rubric and verifier design and cannot establish independent
construct validity.
The primary rater is same-family, sees model identity, and covers only four of
six systems.
The secondary rater is a compact local model, uses a different unit of analysis,
and covers 30 responses from 21 unique cases. Its low agreement cannot be
attributed specifically to family bias. The family-stratified author audit
provides that attribution on identical units, but only for 60 claims and with
author reference scores. No rubric scores exist for the two GPT endpoints, so
six-system conclusions rest on the verifier alone. A same-unit rescoring of all
atomic claims from all six systems by a strong non-Anthropic judge and
independent human ratings are the appropriate next tests.
The structured primary and secondary rating files are preserved, but raw
primary-rater API objects and per-run secondary-rater transcripts were not;
their score aggregation is reproducible while historical execution provenance
is weaker than for the 60 OpenAI extension responses.

The verifier covers four narrow error patterns. Its 160 fixtures exercise known
boundaries, including zero-hit C2 and C4, but seeded text is easier than the full
distribution of natural paraphrases. Because the fixtures were written by the
verifier's authors after the rules were specified, a 160/160 pass shows that
the code implements its intended rules. It cannot show that those rules detect
errors the authors did not anticipate. The \checkc{C3} patterns are
hand-configured for six cases seen during development and partly encode the
expected answers. The audit leaves two of nine flags
unresolved, identifies one definite miss, and leaves one of 24 sampled nonflags
disputed between authors. The
sample is too small for a stable recall estimate. The verifier was developed after the
initial model-rated analysis. A stronger evaluation would freeze version 1 and
apply it without modification to newly authored cases and newly generated
responses.

The 30 cases are synthetic and small; several rubric dimensions exhibit ceiling
effects. The Opus row mixes two served versions. The exploratory
frontier--compact comparison treats repeated case-checks as independent and
groups heterogeneous providers, scales, and serving conditions. It should not
be interpreted as a general model-class effect. None of the results establishes
that any system is suitable for unsupervised evidentiary use.

\section{Conclusion}
An LLM judge related to the systems it scores can produce a leaderboard that
looks precise but has no independent basis for its ordering. In our Android-forensics case study, a
cross-family diagnostic fails to reproduce the primary measurement cleanly. A
direct audit on identical claims shows that the judge's leniency relative to
blinded authors is concentrated in same-family outputs. Meanwhile, four
deterministic checks produce nine auditable flags across six systems.
A model-blinded two-author audit confirms seven, leaves two unresolved, and
finds one miss in 24 sampled nonflags. All 160 seeded conformance fixtures behave as expected, including
positive controls for two checks with no natural flags. These results establish
implementation behavior, specific natural violations, and a measured
same-family leniency on a small audited sample. They do not establish the
validity of the full rubric. A hybrid benchmark that publishes structured
ground truth, applicability sets, witnesses, fixtures, and versioned verifiers
makes that boundary explicit and lets reviewers recompute part of the result
without trusting another generative model.

\section*{Open Science}
The anonymous artifact is available at
\url{https://anonymous.4open.science/r/conflicted-judge-audit-5C2E}. It contains the 30 case JSON files, prompts, 180 response
records, 830 claim annotations for the four primary-scored systems, the
cross-rater sample, verifier source, statistical audit, 160 seeded conformance
fixtures, the blinded author-review packet, the two anonymized completed
author-review records, derived
agreement results, the family-stratified judge-audit script and output,
coordinator key, execution-provenance audit, SHA-256
manifests, and build source.
All 180 response records identify the requested and served model. The 60 OpenAI
extension records additionally preserve unique provider response identifiers,
completion status, timestamps, token usage, and request fingerprints; the 120
legacy records do not preserve provider request identifiers or token usage. API
secrets and provider credentials are excluded. The repository will remain
available throughout review, and a permanent archival identifier will be added
upon publication.

\section*{LLM Usage Considerations}
Language models are the systems under evaluation and two model families also
serve as raters, as described in the evaluation design. Their outputs are
preserved as data rather than accepted as ground truth. The evaluation contains
180 task responses; the bounded design was chosen to compare complete outputs
across fixed cases while limiting unnecessary API and local compute. The local
7.6B Qwen model used a Q4\_K\_M build on an NVIDIA RTX 4050 with 6~GB VRAM; other
systems were accessed through recorded API endpoints. LLMs were used for
editorial purposes in this manuscript, and all outputs were inspected by the
authors to ensure accuracy and originality. During the model-blinded author
audit, the reviewers consulted ordinary websites and reference material. The
authors report that no AI was used to choose or fill review scores or decisions;
any LLM assistance was limited to editorial phrasing after the judgments. We report that exercise
as an author audit because both reviewers are authors. The authors checked
numerical claims against machine-readable outputs and remain responsible for
the text, code, citations, and conclusions.

\section*{Ethical Considerations}
The benchmark is fully synthetic and uses reserved network identifiers. It does
not contain personal data, malware binaries, or operational targets. Forensic
model output can influence investigations, so the benchmark must not be read as
authorization for automated legal or evidentiary decisions. Human examiners
remain responsible for source verification, alternative explanations, and
chain-of-custody requirements.

\begin{thebibliography}{10}
\bibitem{studiawan2025} H. Studiawan, F. Breitinger, and M. Scanlon, ``Towards a standardized methodology and dataset for evaluating LLM-based digital forensic timeline analysis,'' \emph{Forensic Science International: Digital Investigation}, vol. 54, suppl., Art. no. 301982, 2025, doi: 10.1016/j.fsidi.2025.301982.
\bibitem{autodfbench} A. Wickramasekara, A. Densmore, F. Breitinger, H. Studiawan, and M. Scanlon, ``AutoDFBench: A framework for AI generated digital forensic code and tool testing and evaluation,'' in \emph{Proc. Digital Forensics Doctoral Symposium}, 2025, doi: 10.1145/3712716.3712718.
\bibitem{factscore} S. Min \emph{et al.}, ``FActScore: Fine-grained atomic evaluation of factual precision in long form text generation,'' in \emph{Proc. Conference on Empirical Methods in Natural Language Processing}, 2023, pp. 12076--12100, doi: 10.18653/v1/2023.emnlp-main.741.
\bibitem{safe} J. Wei \emph{et al.}, ``Long-form factuality in large language models,'' in \emph{Advances in Neural Information Processing Systems}, vol. 37, 2024, doi: 10.52202/079017-2567.
\bibitem{alce} T. Gao, H. Yen, J. Yu, and D. Chen, ``Enabling large language models to generate text with citations,'' in \emph{Proc. Conference on Empirical Methods in Natural Language Processing}, 2023, pp. 6465--6488, doi: 10.18653/v1/2023.emnlp-main.398.
\bibitem{attrscore} X. Yue, B. Wang, Z. Chen, K. Zhang, Y. Su, and H. Sun, ``Automatic evaluation of attribution by large language models,'' in \emph{Findings of the Association for Computational Linguistics: EMNLP}, 2023, pp. 4615--4635, doi: 10.18653/v1/2023.findings-emnlp.307.
\bibitem{panickssery2024} A. Panickssery, S. R. Bowman, and S. Feng, ``LLM evaluators recognize and favor their own generations,'' in \emph{Advances in Neural Information Processing Systems}, vol. 37, 2024, doi: 10.52202/079017-2197.
\bibitem{zheng2023} L. Zheng \emph{et al.}, ``Judging LLM-as-a-judge with MT-Bench and Chatbot Arena,'' in \emph{Advances in Neural Information Processing Systems}, vol. 36, pp. 46595--46623, 2023, doi: 10.52202/075280-2020.
\bibitem{faireval} P. Wang \emph{et al.}, ``Large language models are not fair evaluators,'' in \emph{Proc. 62nd Annual Meeting of the Association for Computational Linguistics}, 2024, pp. 9440--9450, doi: 10.18653/v1/2024.acl-long.511.
\bibitem{geval} Y. Liu, D. Iter, Y. Xu, S. Wang, R. Xu, and C. Zhu, ``G-Eval: NLG evaluation using GPT-4 with better human alignment,'' in \emph{Proc. Conference on Empirical Methods in Natural Language Processing}, 2023, pp. 2511--2522, doi: 10.18653/v1/2023.emnlp-main.153.
\bibitem{prometheus} S. Kim \emph{et al.}, ``Prometheus: Inducing fine-grained evaluation capability in language models,'' in \emph{Proc. International Conference on Learning Representations}, 2024. [Online]. Available: \url{https://openreview.net/forum?id=8euJaTveKw}.
\bibitem{checklist} M. T. Ribeiro, T. Wu, C. Guestrin, and S. Singh, ``Beyond accuracy: Behavioral testing of NLP models with CheckList,'' in \emph{Proc. Annual Meeting of the Association for Computational Linguistics}, 2020, pp. 4902--4912, doi: 10.18653/v1/2020.acl-main.442.
\bibitem{krippendorff} K. Krippendorff, \emph{Content Analysis: An Introduction to Its Methodology}, 2nd ed. Thousand Oaks, CA, USA: Sage, 2004.
\bibitem{androidsigning} Android Developers, ``Sign your app.'' [Online]. Available: \url{https://developer.android.com/studio/publish/app-signing}. Accessed: Sep. 27, 2026.
\end{thebibliography}

\end{document}
